% Overriding the theorem counter for this chapter to match the 16.1, 16.2 convention

\setcounter{chapter}{15} % Sets chapter counter so the next chapter is 16

\chapter{Linear Maps and Bases}

\subsection{Writing Linear Maps Using Bases and Coordinates}

\renewcommand{\thetheorem}{\thechapter.\arabic{theorem}}

In this lecture, we will assume all vector spaces to be \qt{finite-dimensional} unless otherwise stated. We will write bases of a vector space $V$ as an ordered list (or tuple) $\mathcal{B} = (v_1, \dots, v_n)$.

% def:16.1
\begin{definition}[Coordinate Vector]
\label{def:coordinate_vector}
Let $V$ be a vector space over $K$, and let $\mathcal{B} = (v_1, \dots, v_n)$ be a basis of $V$. Let $v \in V$. We define the \qt{coordinate-vector} of $v$ with respect to the basis $\mathcal{B}$ as
\[
[v]_{\mathcal{B}} = 
\begin{pmatrix}
a_1 \\
\vdots \\
a_n
\end{pmatrix}
\in K^n,
\]
where $a_1, \dots, a_n \in K$ are the unique scalars in $K$ such that $v = a_1 v_1 + \dots + a_n v_n$.
\end{definition}

\begin{remark}
Note that $[v]_{\mathcal{B}}$ is uniquely defined, since every $v \in V$ has a unique representation as a linear combination of the basis vectors $v_1, \dots, v_n$.
\end{remark}

Define a map $\Phi_{\mathcal{B}}: V \to K^n$ given by:
\[
\Phi_{\mathcal{B}}(v) := [v]_{\mathcal{B}} \in K^n.
\]

% prop:16.2
\begin{proposition}[Isomorphism of the Coordinate Map]
\label{prop:coordinate_map_isomorphism}
The map $\Phi_{\mathcal{B}}$ is an isomorphism.
\end{proposition}

\begin{proof}
We first show that $\Phi_{\mathcal{B}}$ is a linear map. Let $a, b \in K$, and let $v, v' \in V$. Write $v$ and $v'$ as linear combinations of the elements of $\mathcal{B}$:
\[
v = a_1 v_1 + \dots + a_n v_n, \quad \text{and} \quad v' = b_1 v_1 + \dots + b_n v_n,
\]
with $a_i, b_i \in K$. By definition, we have:
\[
[v]_{\mathcal{B}} = 
\begin{pmatrix}
a_1 \\
\vdots \\
a_n
\end{pmatrix},
\quad \text{and} \quad
[v']_{\mathcal{B}} = 
\begin{pmatrix}
b_1 \\
\vdots \\
b_n
\end{pmatrix}.
\]
Now, $a v + b v' = (a a_1 + b b_1) v_1 + \dots + (a a_n + b b_n) v_n$, and hence:
\[
\Phi_{\mathcal{B}}(a v + b v') = [a v + b v']_{\mathcal{B}} = 
\begin{pmatrix}
a a_1 + b b_1 \\
\vdots \\
a a_n + b b_n
\end{pmatrix}
= a 
\begin{pmatrix}
a_1 \\
\vdots \\
a_n
\end{pmatrix}
+ b
\begin{pmatrix}
b_1 \\
\vdots \\
b_n
\end{pmatrix}
= a \Phi_{\mathcal{B}}(v) + b \Phi_{\mathcal{B}}(v').
\]
This proves that $\Phi_{\mathcal{B}}$ is linear. 

We now claim that $\Phi_{\mathcal{B}}$ is an isomorphism. By previous results, it is enough to show $\Phi_{\mathcal{B}}$ is bijective. Indeed, if $v \in \ker(\Phi_{\mathcal{B}})$, then $v = 0 v_1 + \dots + 0 v_n = 0$, which implies $\ker(\Phi_{\mathcal{B}}) = \{0\}$. This shows that $\Phi_{\mathcal{B}}$ is injective. 

For the surjectivity of $\Phi_{\mathcal{B}}$, let $(a_1, \dots, a_n)^T \in K^n$. Put $v = a_1 v_1 + \dots + a_n v_n \in V$. Then, $\Phi_{\mathcal{B}}(v) = (a_1, \dots, a_n)^T$, which implies $\operatorname{Im}(\Phi_{\mathcal{B}}) = K^n$.
\end{proof}

\begin{remark}
The inverse map to $\Phi_{\mathcal{B}}$ is given by:
\[
\Phi_{\mathcal{B}}^{-1}: K^n \to V, \quad \Phi_{\mathcal{B}}^{-1}
\begin{pmatrix}
a_1 \\
\vdots \\
a_n
\end{pmatrix}
:= a_1 v_1 + \dots + a_n v_n.
\]
\end{remark}

\begin{remark}
The previous Proposition \ref{prop:coordinate_map_isomorphism} gives yet another proof for the following Corollary.
\end{remark}

% cor:16.3
\begin{corollary}[Classification of Finite-Dimensional Vector Spaces]
\label{cor:vs_isomorphic_to_kn}
Let $V$ be a vector space over $K$ with $\dim V = n \in \mathbb{Z}_{\ge 0}$. Then, $V \cong K^n$.
\end{corollary}

\begin{proof}
Choose a basis $\mathcal{B}$ for $V$. Then, $\Phi_{\mathcal{B}}: V \to K^n$ is an isomorphism.
\end{proof}

\begin{remark}
\begin{enumerate}
    \item \textbf{(a)} The Corollary does NOT hold for $\dim V = \infty$, at least not as stated.
    \item \textbf{(b)} We have seen that $V \cong K^n$ where $n = \dim V < \infty$, but in general, there does not exist a \qt{canonical} (or preferred) isomorphism $V \to K^n$. For example, $\Phi_{\mathcal{B}}$ depends on a choice of a basis $\mathcal{B}$.
\end{enumerate}
\end{remark}

% ex:16.a
\begin{example}[Coordinate Vector Calculation]
\label{ex:coordinate_vector_calculation}
Consider $V = \mathbb{R}^2$, let $\mathcal{E} = \left((1, 0)^T, (0, 1)^T\right)$ be the standard basis, and let $\mathcal{B} = \left((1, 1)^T, (1, -1)^T\right)$ be another basis (check that $\mathcal{B}$ is indeed a basis!). We have
\[
\left[\begin{pmatrix} x \\ y \end{pmatrix}\right]_{\mathcal{E}} = \begin{pmatrix} x \\ y \end{pmatrix}, \quad \text{and we want to find} \quad \left[\begin{pmatrix} x \\ y \end{pmatrix}\right]_{\mathcal{B}} = \begin{pmatrix} a \\ b \end{pmatrix} = ?
\]
Let us first try to determine $[(1, 0)^T]_{\mathcal{B}}$ and $[(0, 1)^T]_{\mathcal{B}}$. We set
\[
\left[\begin{pmatrix} 1 \\ 0 \end{pmatrix}\right]_{\mathcal{B}} = \begin{pmatrix} a_1 \\ a_2 \end{pmatrix}, \quad \text{where} \quad a_1 \begin{pmatrix} 1 \\ 1 \end{pmatrix} + a_2 \begin{pmatrix} 1 \\ -1 \end{pmatrix} = \begin{pmatrix} 1 \\ 0 \end{pmatrix}.
\]
This leads to the system of linear equations:
\[
\begin{cases}
a_1 + a_2 = 1 \\
a_1 - a_2 = 0
\end{cases}
\]
Solving this system yields $a_1 = a_2 = \frac{1}{2}$, and therefore, $[(1, 0)^T]_{\mathcal{B}} = (1/2, 1/2)^T$.

Similarly, for the second standard basis vector, we set
\[
\left[\begin{pmatrix} 0 \\ 1 \end{pmatrix}\right]_{\mathcal{B}} = \begin{pmatrix} b_1 \\ b_2 \end{pmatrix}, \quad \text{where} \quad b_1 \begin{pmatrix} 1 \\ 1 \end{pmatrix} + b_2 \begin{pmatrix} 1 \\ -1 \end{pmatrix} = \begin{pmatrix} 0 \\ 1 \end{pmatrix}.
\]
The resulting system is:
\[
\begin{cases}
b_1 + b_2 = 0 \\
b_1 - b_2 = 1
\end{cases}
\]
Solving this yields $b_1 = \frac{1}{2}$ and $b_2 = -\frac{1}{2}$, so $[(0, 1)^T]_{\mathcal{B}} = (1/2, -1/2)^T$.

Finally, using the fact that $\Phi_{\mathcal{B}}$ is linear, we calculate:
\[
\left[\begin{pmatrix} x \\ y \end{pmatrix}\right]_{\mathcal{B}} = x \left[\begin{pmatrix} 1 \\ 0 \end{pmatrix}\right]_{\mathcal{B}} + y \left[\begin{pmatrix} 0 \\ 1 \end{pmatrix}\right]_{\mathcal{B}} = \begin{pmatrix} x/2 \\ x/2 \end{pmatrix} + \begin{pmatrix} y/2 \\ -y/2 \end{pmatrix} = \begin{pmatrix} \frac{x+y}{2} \\ \frac{x-y}{2} \end{pmatrix}.
\]
\end{example}

\subsection{Representation Matrices}

Now we arrive at the main question. Let $V$ and $W$ be vector spaces over $K$, let $n := \dim V$ and $m := \dim W$, and let $T: V \to W$ be a linear map. Fix a basis $\mathcal{B}$ for $V$ and a basis $\mathcal{C}$ for $W$. Consider the following commutative diagram:
\[
\begin{tikzcd}[row sep=large, column sep=large]
V \arrow[r, "T"] \arrow[d, "\Phi_{\mathcal{B}}"'] & W \arrow[d, "\Phi_{\mathcal{C}}"] \\
K^n \arrow[r, "\Phi_{\mathcal{C}} \circ T \circ \Phi_{\mathcal{B}}^{-1}"'] & K^m
\end{tikzcd}
\]
Since $\Phi_{\mathcal{B}}^{-1}$, $T$, and $\Phi_{\mathcal{C}}$ are all linear maps, the map $\Phi_{\mathcal{C}} \circ T \circ \Phi_{\mathcal{B}}^{-1}: K^n \to K^m$ is also linear. By previous results, any linear map $K^n \to K^m$ can be written using a matrix $A \in \M_{m \times n}(K)$. In other words, there exists a unique matrix $A \in \M_{m \times n}(K)$ such that $\Phi_{\mathcal{C}} \circ T \circ \Phi_{\mathcal{B}}^{-1} = T_A$.

The matrix $A$ is called the \qt{representation matrix} of $T$ with respect to the bases $\mathcal{B}$ and $\mathcal{C}$. We also say that $A$ represents $T$ in the bases $\mathcal{B}$ and $\mathcal{C}$. We denote this matrix by $[T]_{\mathcal{C}}^{\mathcal{B}} \in \M_{m \times n}(K)$.

\begin{question}
\begin{enumerate}
    \item \textbf{(a)} How can we calculate the entries of $A$?
    \item \textbf{(b)} How does $A$ depend on different choices of the bases $\mathcal{B}$ and $\mathcal{C}$?
    \item \textbf{(c)} If we have three vector spaces $V, W, U$ with bases $\mathcal{B}, \mathcal{C}, \mathcal{D}$ respectively, and $T: V \to W$ and $S: W \to U$ are linear maps where $T$ is represented by the matrix $M$ and $S$ by the matrix $N$, then what is the matrix representing $S \circ T: V \to U$ with respect to the bases $\mathcal{B}$ and $\mathcal{D}$?
\end{enumerate}
\end{question}

Let $T: V \to W$ be a linear map, let $\mathcal{B} = (v_1, \dots, v_n)$ be a basis for $V$, and let $\mathcal{C} = (w_1, \dots, w_m)$ be a basis for $W$. For any $v_i \in V$, its image is $T(v_i) \in W$. Under the map $T_A$, the standard basis vector $e_i \in K^n$ maps to $K^m$:
\[
T_A(e_i) = \Phi_{\mathcal{C}} \circ T \circ \Phi_{\mathcal{B}}^{-1}(e_i) = \Phi_{\mathcal{C}}(T(v_i)) = [T(v_i)]_{\mathcal{C}}.
\]
But these are precisely the columns of $A$. In other words:
\[
A = [T]_{\mathcal{C}}^{\mathcal{B}} = 
\begin{pmatrix}
| & & | \\
[T(v_1)]_{\mathcal{C}} & \dots & [T(v_n)]_{\mathcal{C}} \\
| & & |
\end{pmatrix}
\in \M_{m \times n}(K).
\]
Another way to describe this matrix is $A = (a_{ij}) \in \M_{m \times n}(K)$, where the entries $a_{ij}$ are defined uniquely by:
\[
T(v_j) = a_{1j} w_1 + \dots + a_{mj} w_m = \sum_{i=1}^m a_{ij} w_i.
\]

% prop:16.4
\begin{proposition}[Evaluation via Representation Matrix]
\label{prop:matrix_evaluation}
Let $T: V \to W$ be a linear map between two finite-dimensional vector spaces $V$ and $W$. Let $\mathcal{B}$ and $\mathcal{C}$ be bases for $V$ and $W$, respectively. Then, for all $v \in V$:
\[
[T]_{\mathcal{C}}^{\mathcal{B}} \cdot [v]_{\mathcal{B}} = [T(v)]_{\mathcal{C}}.
\]
\end{proposition}

\begin{proof}
The equality $[T]_{\mathcal{C}}^{\mathcal{B}} \cdot [v]_{\mathcal{B}} = [T(v)]_{\mathcal{C}}$ is equivalent to evaluating the linear transformation defined by the matrix $[T]_{\mathcal{C}}^{\mathcal{B}}$ at the vector $[v]_{\mathcal{B}} \in K^n$:
\[
T_{[T]_{\mathcal{C}}^{\mathcal{B}}}([v]_{\mathcal{B}}) = [T(v)]_{\mathcal{C}}.
\]
By definition, $T_{[T]_{\mathcal{C}}^{\mathcal{B}}} = \Phi_{\mathcal{C}} \circ T \circ \Phi_{\mathcal{B}}^{-1}$. But this is precisely how we defined the matrix $[T]_{\mathcal{C}}^{\mathcal{B}}$.
\end{proof}

\begin{example}
\begin{enumerate}
    \item \textbf{(a)} Let $V = K^n$ and $W = K^m$, and let $T: K^n \to K^m$ be given by $T = T_A$ for $A \in \M_{m \times n}(K)$, so $T(v) = A \cdot v$. Take $\mathcal{B} = (e_1, \dots, e_n)$ and $\mathcal{C} = (e_1, \dots, e_m)$ to be the standard bases. Then, $[T]_{\mathcal{C}}^{\mathcal{B}} = A$.
    \item \textbf{(b)} Let $V = K[x]_3$ and $W = K[x]_2$, and let $D: V \to W$ be the derivative map defined by $D(p(x)) = p'(x)$. Let $\mathcal{B} = (1, x, x^2, x^3)$ and $\mathcal{C} = (1, x, x^2)$. We calculate:
    \begin{align*}
    D(1) &= 0 = 0 \cdot 1 + 0 \cdot x + 0 \cdot x^2 \\
    D(x) &= 1 = 1 \cdot 1 + 0 \cdot x + 0 \cdot x^2 \\
    D(x^2) &= 2x = 0 \cdot 1 + 2 \cdot x + 0 \cdot x^2 \\
    D(x^3) &= 3x^2 = 0 \cdot 1 + 0 \cdot x + 3 \cdot x^2
    \end{align*}
    Therefore, the representation matrix is:
    \[
    [D]_{\mathcal{C}}^{\mathcal{B}} = 
    \begin{pmatrix}
    0 & 1 & 0 & 0 \\
    0 & 0 & 2 & 0 \\
    0 & 0 & 0 & 3
    \end{pmatrix}.
    \]
    \setcounter{enumi}{1} % Matching professor numbering
    \item \textbf{(c)} Suppose we view $D$ as a linear map $D: K[x]_3 \to K[x]_3$. Then,
    \[
    [D]_{\mathcal{B}}^{\mathcal{B}} = 
    \begin{pmatrix}
    0 & 1 & 0 & 0 \\
    0 & 0 & 2 & 0 \\
    0 & 0 & 0 & 3 \\
    0 & 0 & 0 & 0
    \end{pmatrix}.
    \]
\end{enumerate}
\end{example}


\subsection{Representing Composition of Linear Maps by Matrices}

Let $T: V \to W$ and $S: W \to U$ be linear maps. Let $\mathcal{A} = (v_1, \dots, v_n)$ be a basis for $V$, let $\mathcal{B} = (w_1, \dots, w_m)$ be a basis for $W$, and let $\mathcal{C} = (u_1, \dots, u_p)$ be a basis for $U$. The dimensions are thus $n = \dim V$, $m = \dim W$, and $p = \dim U$.

Consider the composition $S \circ T: V \longrightarrow U$.

\begin{question}
What is the matrix $[S \circ T]_{\mathcal{C}}^{\mathcal{A}}$? How is it related to the matrices $[T]_{\mathcal{B}}^{\mathcal{A}}$ and $[S]_{\mathcal{C}}^{\mathcal{B}}$?
\end{question}

To answer this, let us write the representation matrices as follows:
\[
[T]_{\mathcal{B}}^{\mathcal{A}} = (a_{ij}), \quad [S]_{\mathcal{C}}^{\mathcal{B}} = (b_{ij}), \quad [S \circ T]_{\mathcal{C}}^{\mathcal{A}} = (c_{ij}).
\]
By definition, the images of the basis vectors are given by:
\begin{equation}
\label{eq:T_basis_image}
T(v_j) = \sum_{k=1}^m a_{kj} w_k, \quad \forall \, 1 \le j \le n,
\end{equation}
\begin{equation}
\label{eq:S_basis_image}
S(w_k) = \sum_{i=1}^p b_{ik} u_i, \quad \forall \, 1 \le k \le m,
\end{equation}
\begin{equation}
\label{eq:composition_basis_image}
(S \circ T)(v_j) = \sum_{i=1}^p c_{ij} u_i, \quad \forall \, 1 \le j \le n.
\end{equation}

But we can also calculate $(S \circ T)(v_j)$ using \eqref{eq:T_basis_image} and \eqref{eq:S_basis_image}:
\begin{align*}
(S \circ T)(v_j) &= S \left( \sum_{k=1}^m a_{kj} w_k \right) \\
&= \sum_{k=1}^m a_{kj} S(w_k) \\
&= \sum_{k=1}^m a_{kj} \left( \sum_{i=1}^p b_{ik} u_i \right) \\
&= \sum_{k=1}^m \sum_{i=1}^p a_{kj} b_{ik} u_i = \sum_{i=1}^p \left( \sum_{k=1}^m b_{ik} a_{kj} \right) \cdot u_i.
\end{align*}
Comparing this with \eqref{eq:composition_basis_image}, we find that $c_{ij} = \sum_{k=1}^m b_{ik} a_{kj}$. In other words, row $i$ of $[S]_{\mathcal{C}}^{\mathcal{B}}$ multiplied by column $j$ of $[T]_{\mathcal{B}}^{\mathcal{A}}$ yields the entry $c_{ij}$:
\[
\begin{pmatrix}
b_{i1} & \dots & b_{im}
\end{pmatrix}
\cdot 
\begin{pmatrix}
a_{1j} \\
\vdots \\
a_{mj}
\end{pmatrix}
= c_{ij}.
\]
This reveals that:
\[
[S]_{\mathcal{C}}^{\mathcal{B}} \cdot [T]_{\mathcal{B}}^{\mathcal{A}} = [S \circ T]_{\mathcal{C}}^{\mathcal{A}},
\]
where the matrix dimensions follow $(p \times m) \cdot (m \times n) = (p \times n)$.

In order to obtain $[S \circ T]_{\mathcal{C}}^{\mathcal{A}}$ from $[T]_{\mathcal{B}}^{\mathcal{A}}$ and $[S]_{\mathcal{C}}^{\mathcal{B}}$, we need to use a new operation on matrices. It is called \qt{matrix multiplication}. But is this operation really new?
- Yes, but not entirely...

